% ============================================================================
% graficas/donut.tex
% Gráfica Donut: Media SUS en grande, arco proporcional, nota y adjetivo
% Puede incluirse con % ============================================================================
% graficas/donut.tex
% Gráfica Donut: Media SUS en grande, arco proporcional, nota y adjetivo
% Puede incluirse con % ============================================================================
% graficas/donut.tex
% Gráfica Donut: Media SUS en grande, arco proporcional, nota y adjetivo
% Puede incluirse con % ============================================================================
% graficas/donut.tex
% Gráfica Donut: Media SUS en grande, arco proporcional, nota y adjetivo
% Puede incluirse con \input{graficas/donut.tex} en documentos y presentaciones
% ============================================================================

\input{graficas/colores.tex}

\begin{tikzpicture}[scale=0.9, every node/.style={transform shape}]
  % Cálculo de ángulo para el arco proporcional (0 a 360 grados)
  \edef\mediaRaw{\directlua{sus.get_stat_raw("media")}}
  \pgfmathsetmacro{\anguloMedia}{\mediaRaw * 3.6}
  \pgfmathsetmacro{\anguloFin}{90 - \anguloMedia}

  % Selección dinámica del color según aceptabilidad
  \pgfmathsetmacro{\colorDonut}{
    \mediaRaw >= 70 ? "susVerde" : (\mediaRaw >= 50 ? "susAmarillo" : "susRojo")
  }

  % Anillo de fondo (100% de la escala)
  \draw[line width=12mm, susGrisFondo] (0,0) circle (2.2cm);

  % Arco de puntuación obtenida
  \ifdim\anguloMedia pt>0.5pt
    \draw[line width=12mm, \colorDonut] (90:2.2cm) arc (90:\anguloFin:2.2cm);
  \fi

  % Textos centrales
  \node[align=center] at (0, 0.25) {
    {\fontsize{26}{30}\selectfont\bfseries \directlua{sus.get_stat("media", 1)}}\\[2pt]
    {\footnotesize\color{gray} sobre 100}
  };

  % Calificación y adjetivo Bangor / Sauro-Lewis en la parte inferior
  \node[align=center, text=susTexto] at (0, -3.2) {
    {\large\bfseries Nota: \directlua{sus.get_stat("nota_media")}}\\[3pt]
    {\normalsize\itshape \directlua{sus.get_stat("adjetivo_media")}}\\[2pt]
    {\footnotesize\bfseries (\directlua{sus.get_stat("aceptabilidad_media")})}
  };
\end{tikzpicture}
 en documentos y presentaciones
% ============================================================================

% ============================================================================
% graficas/colores.tex
% Definición centralizada de colores para las gráficas SUS
% Permite modificar la apariencia de todas las gráficas desde un único lugar
% ============================================================================

\usepackage{xcolor}

% Paleta de aceptabilidad y rendimiento SUS
\definecolor{susRojo}{HTML}{D9534F}       % No aceptable (0 - 50)
\definecolor{susAmarillo}{HTML}{F0AD4E}   % Marginal (50 - 70)
\definecolor{susVerde}{HTML}{5CB85C}      % Aceptable (70 - 100)

% Colores de interfaz y visualización general
\definecolor{susPrincipal}{HTML}{1F4E79}  % Azul corporativo / barras principales
\definecolor{susSecundario}{HTML}{2E75B6} % Azul medio para detalles
\definecolor{susGrisFondo}{HTML}{E9ECEF}  % Gris claro para fondo de escalas / donut
\definecolor{susReferencia}{HTML}{C00000} % Rojo distintivo para línea de media SUS (68)
\definecolor{susTexto}{HTML}{212529}      % Texto oscuro de alto contraste


\begin{tikzpicture}[scale=0.9, every node/.style={transform shape}]
  % Cálculo de ángulo para el arco proporcional (0 a 360 grados)
  \edef\mediaRaw{\directlua{sus.get_stat_raw("media")}}
  \pgfmathsetmacro{\anguloMedia}{\mediaRaw * 3.6}
  \pgfmathsetmacro{\anguloFin}{90 - \anguloMedia}

  % Selección dinámica del color según aceptabilidad
  \pgfmathsetmacro{\colorDonut}{
    \mediaRaw >= 70 ? "susVerde" : (\mediaRaw >= 50 ? "susAmarillo" : "susRojo")
  }

  % Anillo de fondo (100% de la escala)
  \draw[line width=12mm, susGrisFondo] (0,0) circle (2.2cm);

  % Arco de puntuación obtenida
  \ifdim\anguloMedia pt>0.5pt
    \draw[line width=12mm, \colorDonut] (90:2.2cm) arc (90:\anguloFin:2.2cm);
  \fi

  % Textos centrales
  \node[align=center] at (0, 0.25) {
    {\fontsize{26}{30}\selectfont\bfseries \directlua{sus.get_stat("media", 1)}}\\[2pt]
    {\footnotesize\color{gray} sobre 100}
  };

  % Calificación y adjetivo Bangor / Sauro-Lewis en la parte inferior
  \node[align=center, text=susTexto] at (0, -3.2) {
    {\large\bfseries Nota: \directlua{sus.get_stat("nota_media")}}\\[3pt]
    {\normalsize\itshape \directlua{sus.get_stat("adjetivo_media")}}\\[2pt]
    {\footnotesize\bfseries (\directlua{sus.get_stat("aceptabilidad_media")})}
  };
\end{tikzpicture}
 en documentos y presentaciones
% ============================================================================

% ============================================================================
% graficas/colores.tex
% Definición centralizada de colores para las gráficas SUS
% Permite modificar la apariencia de todas las gráficas desde un único lugar
% ============================================================================

\usepackage{xcolor}

% Paleta de aceptabilidad y rendimiento SUS
\definecolor{susRojo}{HTML}{D9534F}       % No aceptable (0 - 50)
\definecolor{susAmarillo}{HTML}{F0AD4E}   % Marginal (50 - 70)
\definecolor{susVerde}{HTML}{5CB85C}      % Aceptable (70 - 100)

% Colores de interfaz y visualización general
\definecolor{susPrincipal}{HTML}{1F4E79}  % Azul corporativo / barras principales
\definecolor{susSecundario}{HTML}{2E75B6} % Azul medio para detalles
\definecolor{susGrisFondo}{HTML}{E9ECEF}  % Gris claro para fondo de escalas / donut
\definecolor{susReferencia}{HTML}{C00000} % Rojo distintivo para línea de media SUS (68)
\definecolor{susTexto}{HTML}{212529}      % Texto oscuro de alto contraste


\begin{tikzpicture}[scale=0.9, every node/.style={transform shape}]
  % Cálculo de ángulo para el arco proporcional (0 a 360 grados)
  \edef\mediaRaw{\directlua{sus.get_stat_raw("media")}}
  \pgfmathsetmacro{\anguloMedia}{\mediaRaw * 3.6}
  \pgfmathsetmacro{\anguloFin}{90 - \anguloMedia}

  % Selección dinámica del color según aceptabilidad
  \pgfmathsetmacro{\colorDonut}{
    \mediaRaw >= 70 ? "susVerde" : (\mediaRaw >= 50 ? "susAmarillo" : "susRojo")
  }

  % Anillo de fondo (100% de la escala)
  \draw[line width=12mm, susGrisFondo] (0,0) circle (2.2cm);

  % Arco de puntuación obtenida
  \ifdim\anguloMedia pt>0.5pt
    \draw[line width=12mm, \colorDonut] (90:2.2cm) arc (90:\anguloFin:2.2cm);
  \fi

  % Textos centrales
  \node[align=center] at (0, 0.25) {
    {\fontsize{26}{30}\selectfont\bfseries \directlua{sus.get_stat("media", 1)}}\\[2pt]
    {\footnotesize\color{gray} sobre 100}
  };

  % Calificación y adjetivo Bangor / Sauro-Lewis en la parte inferior
  \node[align=center, text=susTexto] at (0, -3.2) {
    {\large\bfseries Nota: \directlua{sus.get_stat("nota_media")}}\\[3pt]
    {\normalsize\itshape \directlua{sus.get_stat("adjetivo_media")}}\\[2pt]
    {\footnotesize\bfseries (\directlua{sus.get_stat("aceptabilidad_media")})}
  };
\end{tikzpicture}
 en documentos y presentaciones
% ============================================================================

% ============================================================================
% graficas/colores.tex
% Definición centralizada de colores para las gráficas SUS
% Permite modificar la apariencia de todas las gráficas desde un único lugar
% ============================================================================

\usepackage{xcolor}

% Paleta de aceptabilidad y rendimiento SUS
\definecolor{susRojo}{HTML}{D9534F}       % No aceptable (0 - 50)
\definecolor{susAmarillo}{HTML}{F0AD4E}   % Marginal (50 - 70)
\definecolor{susVerde}{HTML}{5CB85C}      % Aceptable (70 - 100)

% Colores de interfaz y visualización general
\definecolor{susPrincipal}{HTML}{1F4E79}  % Azul corporativo / barras principales
\definecolor{susSecundario}{HTML}{2E75B6} % Azul medio para detalles
\definecolor{susGrisFondo}{HTML}{E9ECEF}  % Gris claro para fondo de escalas / donut
\definecolor{susReferencia}{HTML}{C00000} % Rojo distintivo para línea de media SUS (68)
\definecolor{susTexto}{HTML}{212529}      % Texto oscuro de alto contraste


\begin{tikzpicture}[scale=0.9, every node/.style={transform shape}]
  % Cálculo de ángulo para el arco proporcional (0 a 360 grados)
  \edef\mediaRaw{\directlua{sus.get_stat_raw("media")}}
  \pgfmathsetmacro{\anguloMedia}{\mediaRaw * 3.6}
  \pgfmathsetmacro{\anguloFin}{90 - \anguloMedia}

  % Selección dinámica del color según aceptabilidad
  \pgfmathsetmacro{\colorDonut}{
    \mediaRaw >= 70 ? "susVerde" : (\mediaRaw >= 50 ? "susAmarillo" : "susRojo")
  }

  % Anillo de fondo (100% de la escala)
  \draw[line width=12mm, susGrisFondo] (0,0) circle (2.2cm);

  % Arco de puntuación obtenida
  \ifdim\anguloMedia pt>0.5pt
    \draw[line width=12mm, \colorDonut] (90:2.2cm) arc (90:\anguloFin:2.2cm);
  \fi

  % Textos centrales
  \node[align=center] at (0, 0.25) {
    {\fontsize{26}{30}\selectfont\bfseries \directlua{sus.get_stat("media", 1)}}\\[2pt]
    {\footnotesize\color{gray} sobre 100}
  };

  % Calificación y adjetivo Bangor / Sauro-Lewis en la parte inferior
  \node[align=center, text=susTexto] at (0, -3.2) {
    {\large\bfseries Nota: \directlua{sus.get_stat("nota_media")}}\\[3pt]
    {\normalsize\itshape \directlua{sus.get_stat("adjetivo_media")}}\\[2pt]
    {\footnotesize\bfseries (\directlua{sus.get_stat("aceptabilidad_media")})}
  };
\end{tikzpicture}
